\begingroup
\setlength{\tabcolsep}{4pt}
\begin{tabularx}{\textwidth}{L{0.22\textwidth}X X L{0.18\textwidth}}
\toprule
Módszer & Kulcsfeltétel & Miért lehet gyors? & Időigény \\
\midrule
Leszámláló rendezés & A kulcs egész szám, és $1..k$ közötti tartományba esik. & Nem elempárokat hasonlít, hanem előfordulásszámokat és prefixösszegeket használ. & $\Theta(n+k)$ \\
Számjegyes rendezés & A kulcs $d$ darab pozícióra vagy számjegyre bontható, pozíciónként $k$ lehetséges értékkel. & Több stabil, egyszerű részrendezést végez egymás után. & $\Theta(d(n+k))$ \\
Edényrendezés & A bemenet eloszlásáról is feltételezünk valamit, például egyenletes eloszlást. & Az elemeket edényekbe szórja, és az edényeken belül kicsi rendezéseket végez. & várható $\Theta(n)$ \\
\bottomrule
\end{tabularx}
\endgroup
